\begin{abstract}
Long-running language-model agents may reuse experience after operating
conditions change. Change detection alone, however, cannot determine whether a
stored strategy revision remains useful when its consequences are observed only
after execution. We introduce ShiftMem, a conditional memory lifecycle that uses
public change signals to trigger review, an exact quoted-lead-time condition to
gate retrieval, and delayed operational outcomes to update confidence and
lifecycle state. The LLM is restricted to bounded strategy revisions, while a
shared deterministic controller retains responsibility for daily actions. We
evaluated ShiftMem against an implemented lexical retrieval baseline in a paired
study of a synthetic single-item inventory system frozen before access to Test
outcomes. The evaluation spanned two provider-hosted model deployments, eight
held-out scenarios, and 160 completed deployment--method--scenario--seed runs.
Across 70 model-specific pairs nested in 35 shared environment clusters, the
prespecified but dependence-limited decision rule did
not support a ShiftMem advantage. The mean cost-gap difference (ShiftMem minus
baseline) was $+45.44$ (negative values favor ShiftMem; nominal 95\% mean CI
$[-2.72,93.60]$); the selected two-sided Wilcoxon test, which does not test the
mean directly, gave $p=0.203$. A post-hoc clustered sensitivity retained the
same unfavorable mean estimate, with a 95\% cluster-bootstrap CI of
$[11.26,79.09]$. Variation across deployments and test regimes remained
descriptive. These findings establish neither equivalence nor the intrinsic
ineffectiveness of the conditional memory lifecycle. Instead, within this
testbed, they underscore that conditional memory benefits should be demonstrated rather than
inferred from architecture alone. The study provides an auditable implementation
and evaluation record for testing memory reuse under delayed feedback while
retaining provider and structured-output failures as part of the evaluated
system.
\end{abstract}
